\documentclass[11pt,a4paper]{article}
\usepackage[utf8]{inputenc}
\usepackage{geometry}
\geometry{margin=1in}
\usepackage{amsmath,amssymb,amsfonts}
\usepackage{graphicx}
\usepackage{booktabs}
\usepackage{hyperref}
\usepackage{fancyhdr}

\hypersetup{
    colorlinks=true,
    linkcolor=blue,
    urlcolor=cyan,
    pdftitle={AI Mini-Projects Theoretical Solutions - Makerere University},
    pdfauthor={Group Assignment Team}
}

\pagestyle{fancy}
\fancyhf{}
\rhead{CS221 Aligned AI Mini-Projects}
\lhead{Makerere University}
\rfoot{Page \thepage}

\title{\textbf{Makerere University --- Department of Computer Science}\\
\Large Artificial Intelligence Mini-Project Assignments\\
\large Theoretical Solutions Master Document}
\author{\textbf{Team Submission}}
\date{Deadline: 5th December 2026}

\begin{document}

\maketitle

\section*{Group Members and Task Allocation}
\begin{table}[h!]
\centering
\small
\begin{tabular}{@{}lllll@{}}
\toprule
\textbf{\#} & \textbf{Full Name} & \textbf{Reg / Student No.} & \textbf{Email} & \textbf{Assigned Module} \\ \midrule
1 & Henry & [Reg No] & [Email] & 1. Foundations \\
2 & [Member Name] & [Reg No] & [Email] & 2. Sentiment Classification \\
3 & Precious Atwiine & [Reg No] & [Email] & 3. Route Planning \\
4 & Marcus & [Reg No] & [Email] & 4. MountainCar \\
5 & [Co-Lead] & [Reg No] & [Email] & 4. MountainCar \\
6 & Jill & [Reg No] & [Email] & 5. Pac-Man \\
7 & Mukisa Percy & [Reg No] & mukisapercy4@gmail.com & 5. Pac-Man \\
8 & Joseph Mary & [Reg No] & [Email] & 6. Course Scheduling / Bayesian \\
9 & Josh Kwizera & [Reg No] & [Email] & 7. Car Tracking / Logic \\
10 & [Integration Lead] & [Reg No] & [Email] & 8. Integration \& Submission Lead \\ \bottomrule
\end{tabular}
\end{table}

\noindent\textbf{Course Portal:} \url{https://muele.mak.ac.ug/course/section.php?id=230921}\\
\textbf{Reference Coursework:} \url{https://stanford-cs221.github.io/spring2026/#coursework}

\newpage
\tableofcontents
\newpage

\section{Assignment 1: Foundations (\texttt{foundations})}
\textit{Assigned to: Henry}
\subsection{Mathematical Optimization}
\subsection{Probability and Maximum Likelihood}

\section{Assignment 2: Sentiment Classification (\texttt{sentiment})}
\textit{Assigned to: [Assigned Member]}
\subsection{Loss Function Derivations}
\subsection{Stochastic Gradient Descent Analysis}

\section{Assignment 3: Route Planning (\texttt{route})}
\textit{Assigned to: Precious Atwiine}
\subsection{Heuristic Admissibility and Consistency Proofs}
\subsection{Search Complexity Comparison}

\section{Assignment 4: Controlling MountainCar (\texttt{mountaincar})}
\textit{Assigned to: Marcus and Co-Lead}
\subsection{Bellman Optimality Equations}
\subsection{Convergence of Value Iteration and Q-Learning}

\section{Assignment 5: Pac-Man Multi-Agent Adversarial Search (\texttt{pacman})}
\textit{Assigned to: Jill and Mukisa Percy}

\subsection{Multi-Agent Minimax Game Tree Formulation}
In the multi-agent Pac-Man environment, games are zero-sum with one maximizing agent (Pac-Man, index $0$) and $k$ minimizing adversaries (Ghosts, indices $1, \dots, k$). One search ply consists of $1 + k$ agent turns.

The value of state $s$ is recursively given by:
\begin{equation}
V(s) = \begin{cases}
\text{Utility}(s), & \text{if } s \text{ is terminal or depth limit } d = d_{\max} \\
\max_{a \in \text{Actions}(s, 0)} V(\text{Succ}(s, a)), & \text{if agent is Pac-Man (index 0)} \\
\min_{a \in \text{Actions}(s, i)} V(\text{Succ}(s, a)), & \text{if agent is ghost } i \in \{1, \dots, k\}
\end{cases}
\end{equation}

\noindent\textbf{Sub-optimal Adversaries:} If ghosts play sub-optimally (e.g., random walks), the Minimax agent guarantees a realized utility at least as great as the minimax theoretical value:
\begin{equation}
\text{Utility}_{\text{actual}} \ge V_{\text{minimax}}(s_0)
\end{equation}
However, Minimax is inherently pessimistic and cannot exploit predictable ghost mistakes, which Expectimax handles naturally.

\subsection{Alpha-Beta Pruning Correctness and Complexity}
Alpha-Beta pruning tracks parameters $\alpha$ (best choice for MAX along the path) and $\beta$ (best choice for MIN along the path):
\begin{itemize}
    \item \textbf{$\beta$-Pruning at MAX Node:} Truncates search if $v \ge \beta$. An ancestor MIN node would never choose this path.
    \item \textbf{$\alpha$-Pruning at MIN Node:} Truncates search if $v \le \alpha$. An ancestor MAX node already found a superior option elsewhere.
\end{itemize}

\noindent Under optimal move ordering (best successors explored first), the effective branching factor is halved:
\begin{equation}
T_{\text{optimal}}(d) = O\left(b^{d/2}\right) = O\left((\sqrt{b})^d\right)
\end{equation}
allowing the search to double in tree depth within identical time constraints.

\subsection{Expectimax vs. Minimax Invariance Under Transformations}
\begin{enumerate}
    \item \textbf{Monotonic Transformations:} Minimax decisions depend solely on value orderings. For any strictly increasing function $f$, $\arg\max_a f(V(s_a)) = \arg\max_a V(s_a)$. Minimax is invariant under all monotonic transformations.
    \item \textbf{Expectimax Sensitivity:} Expectimax calculates expectations:
    \begin{equation}
    \mathbb{E}[V] = \sum_a P(a) V(s_a)
    \end{equation}
    Because $\mathbb{E}[f(X)] \neq f(\mathbb{E}[X])$ for non-linear functions (Jensen's Inequality), Expectimax decisions can change under non-linear monotonic transformations. Expectimax is invariant \textbf{only} under positive affine transformations ($f(x) = c x + d, c > 0$).
\end{enumerate}

\subsection{State Evaluation Function Formulation}
Our heuristic evaluation function combines domain features linearly:
\begin{equation}
\text{Eval}(s) = w_1 \cdot \text{Score}(s) + \frac{w_2}{d_{\text{food}}(s) + 1} - w_3 N_{\text{food}}(s) - w_4 N_{\text{capsules}}(s) + \sum_{g} \text{GhostPenalty}(g, s)
\end{equation}

\section{Assignment 6: Course Scheduling \& Bayesian Networks (\texttt{scheduling} / \texttt{bayesian})}
\textit{Assigned to: Joseph Mary}
\subsection{Constraint Satisfaction Problem Modeling}
\subsection{Bayesian Network Independence and Inference}

\section{Assignment 7: Car Tracking \& Logic (\texttt{tracking} / \texttt{logic})}
\textit{Assigned to: Josh Kwizera}
\subsection{Hidden Markov Models and Particle Filtering}
\subsection{First-Order Logic Transformations and Resolution}

\end{document}
